\documentclass[10pt]{article}

% Packages pour les marges
\usepackage[
    top=2cm,
    bottom=2cm,
    left=2cm,
    right=2cm
]{geometry}

% Police sans serif
% \usepackage{helvet}
% \renewcommand{\familydefault}{\sfdefault}

% Packages existants
\usepackage[french]{babel}
\usepackage[utf8]{inputenc}
\usepackage[T1]{fontenc}
\usepackage{amsmath}
\usepackage{amsfonts}
\usepackage{amssymb}
\usepackage[version=4]{mhchem}
\usepackage{stmaryrd}
\usepackage{listings}

\usepackage{enumitem}
\usepackage{ifthen}
\usepackage{graphicx}
\usepackage{verbatim}
\usepackage{url}
\usepackage{mdframed}
\usepackage[sf]{titlesec}
\titleformat{\section}{\sffamily\large\bfseries}{\thesection}{1em}{}
% Définition de la variable pour afficher les corrections
\newboolean{showSolutions}
% Décommentez la ligne suivante pour afficher les solutions
\input \jobname.adr
% \input \jobname.adr


\title{TD3 - Problèmes à traiter à plusieurs}

\author{}
\date{}

\newenvironment{solution}
    {\par\vspace{0.5em}\begin{mdframed}[linewidth=0.5pt]\par}
    {\end{mdframed}\par\vspace{0.5em}}

\begin{document}
\sffamily
\maketitle






\section*{Problème 1 - Fabriquer des boulons}
Un industriel fabrique des boulons dont une certaine proportion, $5 \%$, est hors des tolérances fixées par le cahier des charges. Pour réceptionner un lot, le client prélève sur ce lot un échantillon de 30 boulons et il refuse le lot si le nombre de boulons défectueux trouvés dans cet échantillon est strictement supérieur à 4.
\begin{enumerate}
    \item Quelle est la probabilité qu'un lot soit refusé?
    \item L'industriel fabrique 1000 lots par mois. Déterminer (en justifiant) la loi du nombre $N$ de lots refusés en un mois.
    \item Sachant que chaque lot refusé entraîne, pour l'industriel, une perte de $100 €$, quelle est la perte moyenne mensuelle supportée par le fabricant, du fait de l'imperfection de sa production?
    \item Quelle est la probabilité que l'industriel supporte une perte d'au moins $500 €$ en un mois du fait de l'imperfection de sa production?
\end{enumerate}
N.B. : On supposera, pour simplifier, que les états respectifs des boulons (dans ou hors les normes du cahier des charges) sont mutuellement indépendants.

\vspace{2cm}

\section*{Problème 2 - Surbooking}
Une ligne aérienne est équipée d'Airbus A320 dotés de 174 places. On a observé qu'en moyenne $3 \%$ des personnes ayant acheté un billet ne se présentent pas à l'embarquement. Pour augmenter son bénéfice, la compagnie aérienne pratique le surbooking: elle vend deux billets de plus qu'il n'y a de places, soit 176 billets pour chaque vol. Pour un vol donné, on note $X_i$ la variable aléatoire valant 1 si le détenteur de la $i^{\mathrm{e}}$ billet ne se présente pas à l'embarquement et 0 sinon.
\begin{enumerate}
    \item En supposant, pour simplifier, les $X_i$ indépendantes, déterminer la loi du nombre $N$ de détenteurs de billets pour un vol ne se présentant pas à l'embarquement.
    \item Quelle est la probabilité pour qu'il manque exactement une place dans l'avion?
    \item Quelle est la probabilité pour qu'il manque exactement deux places dans l'avion?
    \item Quelle est la probabilité pour qu'il y ait assez de place dans l'avion pour toutes les personnes se présentant à l'embarquement?
    \vspace{1em}
    
    Lorsqu'une personne munie d'un billet ne peut embarquer faute de place, elle est indemnisée à hauteur de $x$ euros. Par ailleurs, les billets sont vendus 90 euros l'unité. On note $D$ le montant total des dédommagements versés lors d'un vol sur cette ligne.
    \item Déterminer la loi de $D$.
    \item Exprimer l'espérance de $D$ en fonction de $x$.
    \item Pour quelles valeurs de $x$, la stratégie de surbooking est-elle en moyenne rentable sur cette ligne?
\end{enumerate}

\newpage

\section*{Problème 3 - coût de dépistage}
Le coût de dépistage de la maladie $M$ à l'aide d'un test sanguin est $c$. La probabilité qu'une personne soit atteinte de la maladie $M$ est $p$. Chaque personne est malade indépendamment des autres. Pour effectuer un dépistage parmi $N$ personnes, on propose les deux stratégies suivantes.

Stratégie 1 : Un test par personne.
Stratégie 2: On suppose que $N / n$ est entier et on regroupe les $N$ personnes en $N / n$ groupes de $n$ personnes. Par groupe, on mélange les prélèvements sanguins des $n$ personnes et on effectue le test. Si on détecte la maladie $M$ dans le mélange, alors on refait un test sanguin pour chacune des $n$ personnes. On supposera que les personnes sont infectées ou non indépendamment les unes des autres.
\begin{enumerate}
    \item Calculer le coût de la Stratégie 1.
    \item Déterminer la probabilité pour qu'un groupe de $n$ personnes contienne au moins une personne infectée.
    \item Soit $X$ la variable aléatoire donnant le nombre de groupes pour lequel des tests individuels seront effectués. Déterminer la loi de $X$
    \item En déduire le coût moyen de la Stratégie 2.
    \item On supposera $n p \ll 1$, et on montrera que $n \simeq 1 / \sqrt{p}$ est une taille qui minimise correctement le coût moyen de la stratégie 2.
    Indication: On pourra utiliser (en le justifiant) que si $n p \ll 1$, alors

$$
c N\left(1+\frac{1}{n}-(1-p)^n\right)=c N\left(n p+\frac{1}{n}\right)+o(n p)
$$

    \item Quelle stratégie choisissez-vous? Illustrer vos conclusions pour le cas où $p=5 \%$ en calculant le pourcentage moyen d'économie réalisée.
\end{enumerate}

\vspace{2cm}
\section*{Problème 4 - Sondage et élection}
Lors d'une élection entre deux candidats $A$ et $B$, on réalise un sondage en interrogeant $n$ des $N$ votants à la sortie des urnes. On note $X_n$ le nombre de réponses au sondage favorables au candidat $A$.
\begin{enumerate}
    \item On suppose que les sondés sont sélectionnés indépendamment et avec remise de sorte à répéter $n$ fois indépendamment la même expérience de Bernoulli. Déterminer la loi de $X_n$.
    \item Quel est le comportement asymptotique de la moyenne empirique $f^{(n)}=\frac{X_n}{n}$ des réponses favorables au candidat $A$ ?
    \item Déterminer un intervalle centré en $f^{(n)}$ contenant, avec probabilité asymptotique de 0,95 , la proportion $p$ de votants favorables au candidat $A$.
    
    Note: On pourra utiliser que, si $Z \sim \mathcal{N}(0,1)$, alors $\mathbf{P}[-1,96 \leq Z \leq 1,96] \simeq 0,95$.
    
    Remarque : On montre que si, comme souvent, le sondage est effectué sans remise cet intervalle reste correct, bien que l'indépendance des réponses ne soit pas garantie, dès que $N \gg n$.
    
    \item En déduire la précision à laquelle peut prétendre avec une approximation asymptotique un sondage sur deux candidats effectué, au risque d'erreur de $5 \%$, sur un échantillon de 1000 personnes.
    \item Est-ce que ce résultat dépend de la taille $N$ de la population?
    \item En Arizona, pour l'élection présidentielle américaine de 2020, on comptait 5083137 inscrits et 3329455 votants. En négligeant la présence d'autres candidats et sachant qu'il y a eu 18610 voix d'écart entre les deux candidats principaux, quel est le nombre de personnes $n$ qu'il aurait fallu interroger à la sortie des urnes dans un sondage pour garantir avec au moins $95 \%$ de chance que le candidat in fine vainqueur allait l'être?
\end{enumerate}




\end{document}